\medskip

On compte quatre groupes sanguins dans l'espèce humaine: A, B, AB et O.

Chaque groupe sanguin peut présenter un facteur rhésus. Lorsqu'il est présent, on dit que le rhésus est positif, sinon on dit qu'il est négatif.

Au sein de la population française, on sait que:


\begin{itemize}[label=$\bullet~$]
\item 45\,\% des individus appartiennent au groupe A, et parmi eux 85\,\% sont de rhésus positif;
\item 10\,\% des individus appartiennent au groupe B, et parmi eux 84\,\% sont de rhésus positif;
\item 3\,\% des individus appartiennent au groupe AB, et parmi eux 82\,\% sont de rhésus positif.
\end{itemize}
\smallskip

On choisit au hasard une personne dans la population française. 

On désigne par:

\begin{itemize}[label=$\bullet~$]
\item $A$ l'évènement \og La personne choisie est de groupe sanguin A \fg{} ;
\item $B$ l'évènement \og La personne choisie est de groupe sanguin B \fg ;
\item $AB$ l'évènement \og La personne choisie est de groupe sanguin AB \fg ; 
\item $O$ l'évènement \og La personne choisie est de groupe sanguin O \fg ;
\item $R$ l'évènement \og La personne choisie a un facteur rhésus positif \fg.
\end{itemize}

\medskip

Pour un évènement quelconque $E$, on note $\overline{E}$ l'évènement contraire de $E$ et $p(E)$ la probabilité de $E$.

\medskip

\begin{minipage}{0.55\linewidth}
\begin{enumerate}
\item Recopier l'arbre ci-contre en complétant les dix pointillés.
\item Montrer que $p(B \cap R) = 0,084$. Interpréter ce résultat dans le contexte de l'exercice.
\item On précise que $p(R) = \np{0,8397}$.

Montrer que $p_O(R) = 0,83$.
\item On dit qu'un individu est \og donneur universel \fg lorsque son sang peut être transfusé à toute personne sans risque d'incompatibilité.

Le groupe O de rhésus négatif est le seul vérifiant cette caractéristique.

Montrer que la probabilité qu'un individu choisi au hasard dans la population française soit donneur universel est de \np{0,0714}.
\end{enumerate}
\end{minipage}\hfill
\begin{minipage}{0.42\linewidth}
\begin{center}
\pstree[treemode=R,levelsep=3cm]{\TR{}}
{\pstree{\TR{$A~$}\taput{\ldots}}
	{\TR{$R$} \taput{\ldots}
	\TR{$\overline{R}$} \tbput{\ldots}
	}
\pstree{\TR{$B~$}\taput{\ldots}}
	{\TR{$R$} \taput{\ldots}
	\TR{$\overline{R}$} \tbput{\ldots}
	}
\pstree{\TR{$AB~$}\tbput{\ldots}}
	{\TR{$R$} \taput{\ldots}
	\TR{$\overline{R}$} \tbput{\ldots}
	}
	\pstree{\TR{$O~$}\tbput{\ldots}}
	{\TR{$R$} \taput{}
	\TR{$\overline{R}$} \tbput{}
	}
}
\end{center}
\end{minipage}


\begin{enumerate}[resume]
\item Lors d'une collecte de sang, on choisit un échantillon de $100$ personnes dans la population d'une ville française. Cette population est suffisamment grande pour assimiler ce choix à un tirage avec remise.
On note $X$ la variable aléatoire qui à chaque échantillon de $100$~personnes
associe le nombre de donneurs universels dans cet échantillon.

	\begin{enumerate}
		\item Justifier que $X$ suit une loi binomiale dont on précisera les paramètres.
		\item Déterminer à $10^{-3}$ près la probabilité qu'il y ait au plus 7 donneurs universels
dans cet échantillon.
		\item Montrer que l'espérance $E(X)$ de la variable aléatoire $X$ est égale à $7,14$ et
que sa variance $V(X)$ est égale à $6,63$ à $10^{-2}$ près.
	\end{enumerate}
\item Lors de la semaine nationale du don du sang, une collecte de sang est organisée dans $N$ villes françaises choisies au hasard numérotées 1,2, 3, ... , $N$ où $N$ est un entier naturel non nul.

On considère la variable aléatoire $X_1$ qui à chaque échantillon de 100 personnes de la ville 1 associe le nombre de donneurs universels dans cet échantillon.

On définit de la même manière les variables aléatoires $X_2$ pour la ville 2, ... , $X_N$ pour la ville $N$.

On suppose que ces variables aléatoires sont indépendantes et qu'elles admettent la même espérance égale à $7,14$ et la même variance égale à $6,63$.

On considère la variable aléatoire $M_N = \dfrac{X_1 + X_2 + \ldots + X_N}{N}$.
	\begin{enumerate}
		\item Que représente la variable aléatoire $M_N$ dans le contexte de l'exercice ?
		\item Calculer l'espérance $E(M_N)$.
		\item On désigne par $V(M_N)$ la variance de la variable aléatoire $M_N$.

Montrer que $V(M_N) = \dfrac{6,63}{N}$.
		\item Déterminer la plus petite valeur de $N$ pour laquelle l'inégalité de Bienaymé- Tchebychev permet d'affirmer que:

\[P(7 < M_N < 7,28) \geqslant 0,95.\]

	\end{enumerate}
\end{enumerate}

\bigskip
